\section{Typing rules and safety}\label{sec:05-rigor-check:03-typing-rules-and-safety:1}

Chapters 1--4 have relied on Lean's type checking constantly without ever seeing its rules written down. This section makes two things precise: the actual rules Lean's kernel checks a term against (using a small, representative fragment, the \textbf{simply typed λ-calculus}, STLC), and the specific rule governing the universe hierarchy §2 just introduced informally.

An untyped calculus has a serious defect: it permits terms that reduce forever (Chapter 1 §4 mentioned β-reduction; nothing there stopped a term from β-reducing to itself, endlessly), and nothing prevents applying one kind of value where a completely different kind was intended. STLC fixes this by attaching a \textbf{type} to every term and requiring application to respect types --- the direct ancestor of Chapter 1's ``every expression has a type.''

\subsection{Types}\label{sec:05-rigor-check:03-typing-rules-and-safety:2}

\[
\tau ::= \iota \;\mid\; \tau \to \tau
\]

where \(\iota\) ranges over some fixed collection of \textbf{base types} (think \texttt{Nat}, \texttt{Bool}) and \(\tau_1 \to \tau_2\) is the type of functions from \(\tau_1\) to \(\tau_2\) --- exactly Lean's \texttt{→}. Every type is built from base types and arrows. There is not yet any way to quantify over \emph{all} types the way Chapter 1's \texttt{identity \{α : Type\} (x : α) : α := x} did --- that extra generality is exactly Chapter 1 §5's Π-types, already covered.

\begin{progcorner}

Python's own type hints plus \texttt{mypy} are a light version of exactly this system, worth seeing first since this Python tooling runs today, without any Lean installation at all:

\end{progcorner}

\begin{lstlisting}[language=Python,style=python]
def apply_twice(f: int, x: int) -> int:  # pretend f is Callable[[int], int]
    ...

def to_str(n: int) -> str:
    return str(n)

# mypy accepts this: to_str's declared input type (int) matches what
# it is given (an int literal).
to_str(5)

# mypy rejects this with an error, WITHOUT running anything:
# error: Argument 1 to "to_str" has incompatible type "str"; expected "int"
to_str("already a string")
\end{lstlisting}

\texttt{mypy}'s check on \texttt{to\_\allowbreak{}str("already a string")} is exactly the (App) rule below, refusing to combine a function with an argument of the wrong type, caught by reading the code rather than running it. The difference is that Python's hints are optional and only checked when \texttt{mypy} is run at all; Lean's type system is not optional, is checked on every single elaboration, and is a real part of the language's meaning rather than an add-on linter. STLC below is what is actually going on, underneath both.

\subsection{Typing judgments and rules}\label{sec:05-rigor-check:03-typing-rules-and-safety:3}

A \textbf{typing judgment} \(\Gamma \vdash t : \tau\) reads ``in context \(\Gamma\) (a list of variable-type assignments \(x_1:\tau_1, \ldots, x_n:\tau_n\)), the term \(t\) has type \(\tau\).'' This is precisely what Lean's \texttt{\#check} reports, and precisely the ``context'' that appears above the line in every tactic goal state since Chapter 4. The rules that generate valid judgments:

\[
\dfrac{x : \tau \in \Gamma}{\Gamma \vdash x : \tau} \;\text{(Var)}
\qquad
\dfrac{\Gamma, x:\tau_1 \vdash t : \tau_2}{\Gamma \vdash \lambda x.\, t : \tau_1 \to \tau_2} \;\text{(Abs)}
\qquad
\dfrac{\Gamma \vdash t_1 : \tau_1 \to \tau_2 \quad \Gamma \vdash t_2 : \tau_1}{\Gamma \vdash t_1\, t_2 : \tau_2} \;\text{(App)}
\]

Each rule, read as a Lean fact already familiar from earlier chapters:

\begin{itemize}
\tightlist
\item
  \textbf{(Var)} --- if \texttt{x}'s type is recorded in the local context, \texttt{\#check x} reports that type; there is nothing to derive.
\item
  \textbf{(Abs)} --- to type-check \texttt{fun x => t}, extend the context with \texttt{x : τ1} (a fresh assumption, exactly like a hypothesis introduced by \href{https://lean-lang.org/doc/reference/latest/Tactic-Proofs/Tactic-Reference/}{\texttt{intro}} in Chapter 4), check that \texttt{t}'s type is \texttt{τ2} under that extended context, and conclude the whole abstraction has type \texttt{τ1 → τ2}. This \emph{is} how Lean checks every \texttt{def f (x : τ1) : τ2 := t} written since Chapter 1.
\item
  \textbf{(App)} --- to type-check \texttt{t1 t2}, \texttt{t1} must have a function type whose domain matches \texttt{t2}'s type exactly; the result has the codomain type. This is exactly the error Chapter 4 discussed under ``reading a tactic failure'': \href{https://lean-lang.org/doc/reference/latest/Tactic-Proofs/Tactic-Reference/}{\texttt{exact}} \texttt{e} fails with a \textbf{type mismatch} precisely when this rule's side condition (that \(\tau_1\) must match on both sides) is not met.
\end{itemize}

\subsection{Type preservation and progress: the payoff}\label{sec:05-rigor-check:03-typing-rules-and-safety:4}

Two theorems about STLC are the entire reason to bother with a type system at all:

\begin{itemize}
\tightlist
\item
  \textbf{Progress}: a well-typed closed term (no free variables) is either already a \textbf{value} --- an abstraction, or (if base types come with their own constants, as \texttt{Nat}/\texttt{Bool} effectively do) a constant of a base type --- or it can take a β-reduction step. It never ``gets stuck'' partway through evaluation. There is no well-typed analogue of ``apply \texttt{3} to \texttt{true},'' because (App)'s side condition would already have rejected such a term at elaboration time, before any reduction is attempted.
\item
  \textbf{Preservation} (subject reduction): if \(\Gamma \vdash t : \tau\) and \(t \longrightarrow_\beta t'\), then \(\Gamma \vdash t' : \tau\). Reduction never changes a term's type. This is \emph{exactly} why this chapter's own definitional equality (§4, next) is trustworthy: reducing a term to normal form (what \texttt{\#eval}/\texttt{rfl} do) can never accidentally produce something of a different type than the one it started with.
\end{itemize}

Together, progress and preservation are the formal content of ``well-typed programs do not go wrong,'' and, viewed through Curry--Howard (Chapter 3), ``well-typed proof terms do not prove something other than what they claim to prove.'' The entire reliability of Lean as a proof assistant rests on an extended version of exactly these two theorems, proved once for Lean's kernel and then trusted for every proof checked thereafter.

\subsection{What STLC still cannot do}\label{sec:05-rigor-check:03-typing-rules-and-safety:5}

STLC cannot type \texttt{identity} from Chapter 1 \emph{polymorphically}. One could write \texttt{identity\_\allowbreak{}Nat : Nat → Nat} and separately \texttt{identity\_\allowbreak{}Bool : Bool → Bool}, one arrow-type definition per base type. But there is no single term of a single STLC type that captures ``the identity function, at every type.''

\begin{progcorner}

Plain Python never runs into this, because it has no static types to begin with --- \texttt{def identity(x): return x} already works on anything, at runtime, with zero declarations. But the instant type hints are added, wanting \texttt{mypy} to check the \emph{general} claim ``this returns whatever type it was given'' requires a dedicated feature, \texttt{TypeVar}, precisely because bare hints have the same limitation as STLC:

\end{progcorner}

\begin{lstlisting}[language=Python,style=python]
from typing import TypeVar
T = TypeVar("T")

def identity(x: T) -> T:   # one signature, valid at every type
    return x

identity(5)        # T := int
identity("hello")  # T := str
\end{lstlisting}

Without \texttt{TypeVar}, the only option would be writing \texttt{identity\_\allowbreak{}int(x: int) -\allowbreak{}> int} and \texttt{identity\_\allowbreak{}str(x: str) -\allowbreak{}> str} separately --- exactly STLC's one-arrow-type-per-base-type wall. \texttt{TypeVar} is Python typing's escape hatch for this specific gap, and it is a useful anchor: a much lighter version of the same extra generality Lean's \texttt{identity \{α : Type\} (x : α) : α := x} uses, where \texttt{α} is filled in silently every call, the same way \texttt{mypy} silently solves \texttt{T := int} above.

This is precisely the gap \hyperref[sec:01-basics:05-pi-sigma-and-coc:1]{Chapter 1 §5} already closed: \textbf{dependent types} let a type itself depend on a term (here, the type argument \texttt{α}). That is exactly the extra generality \texttt{identity \{α : Type\} (x : α) : α := x} uses, and it is unavailable in STLC (or in Python's \texttt{TypeVar}, which is real but considerably less powerful --- it cannot let a \emph{return type} depend on an ordinary \emph{value} argument the way Chapter 1 §3's \texttt{Vec.replicate} does).

\subsection{Universes, as a typing rule}\label{sec:05-rigor-check:03-typing-rules-and-safety:6}

§2 introduced the hierarchy \(\mathtt{Type} : \mathtt{Type}\,1 :
\mathtt{Type}\,2 : \cdots\) to avoid a Russell-style paradox (a ``type of all types'' that contains itself leads to the same contradiction as ``the set of all sets that do not contain themselves''), and showed \texttt{Group : Type → Type} had to live in \texttt{Type 1}. In the calculus of constructions --- the formal system CoC/CIC that Chapter 1 §5 named, extending STLC above with Π-types and universes --- this is stated as a typing rule for the universes themselves:

\[
\dfrac{}{\mathtt{Type}\,i : \mathtt{Type}\,(i+1)}
\]

together with a rule saying Π-types (function types, including ordinary \texttt{→}) stay inside a suitable universe:

\[
\dfrac{\Gamma \vdash A : \mathtt{Type}\,i \quad \Gamma, x:A \vdash B : \mathtt{Type}\,j}
      {\Gamma \vdash \big(\textstyle\prod_{x:A} B\big) : \mathtt{Type}\,(\max(i,j))}
\]

This is exactly §2's \texttt{Group : Type → Type} computation spelled out as a general rule: with \(A = \mathtt{Type}\) (itself living in \texttt{Type 1}) and \(B = \mathtt{Type}\) again, the rule gives \(\max(1, 1) = 1\), so \texttt{Type → Type} lands in \texttt{Type 1}, one level above \texttt{Type} itself.

\begin{progcorner}

In Python, \texttt{type(3)} is \texttt{int}, and \texttt{type(int)} is \texttt{type} --- and, unlike Lean's stratified hierarchy, \texttt{type(type)} is \emph{also} \texttt{type}:

\end{progcorner}

\begin{lstlisting}[language=Python,style=python]
>>> type(3)
<class 'int'>
>>> type(int)
<class 'type'>
>>> type(type)
<class 'type'>
\end{lstlisting}

Python allows \texttt{type} to be its own type, with no stratification at all, because Python's type system is not being used as a proof system --- there is no soundness property at stake that a Russell-style paradox could break. Lean's \texttt{Type} cannot self-apply this way (\texttt{Type : Type} is \emph{inconsistent} --- it allows encoding Girard's paradox and proving \texttt{False}), which is exactly why the infinite, strictly increasing hierarchy above is load-bearing rather than pedantry. This is one of the few places where the Python comparison genuinely runs out: it is not that Python does the same thing more simply, it is that Python does not need to solve this problem at all, because nothing checks proofs against it.

\subsection{References}\label{sec:05-rigor-check:03-typing-rules-and-safety:7}

Full citations in the \hyperref[sec:root:bibliography:1]{Bibliography}.

\begin{itemize}
\tightlist
\item
  Pierce (\cite{Pierce2002}), Ch. 9--11 --- the standard reference for STLC, including the (Var)/(Abs)/(App) rules and full proofs of progress and preservation, in the exact form used in this section.
\item
  Milner (\cite{Milner1978}) --- the theoretical background for why STLC alone cannot type polymorphic functions like \texttt{identity}.
\item
  Python \texttt{typing} module documentation and mypy documentation (\cite{PythonTyping}, \cite{MypyDocs}) --- for the Python-side comparison used in this section's boxes.
\item
  Coquand and Huet (\cite{CoquandHuet1988}) --- the original paper defining CoC, whose universe-formation rule is stated above.
\item
  Girard (\cite{Girard1971}) --- the source of ``Girard's paradox,'' the inconsistency \texttt{Type : Type} would introduce.
\item
  \emph{Theorem Proving in Lean 4} (\cite{TPIL4}), ``Dependent Types'' --- Lean's own documentation on universes, matching the presentation here.
\end{itemize}
